
\subsection{模型基本假设}
模型遵循题目计费规则，并作如下假设以便复现：

\begin{enumerate}
    \item \textbf{交易与费用}：微网只能购电，不能售电；多余光伏可弃用，未消纳的合同电量仍按原价计费。模型只统计合同购电费和应急购电费。
    \item \textbf{储能参数}：容量为$12000\ \mathrm{kWh}$，运行范围为$[1200,10800]\ \mathrm{kWh}$；充、放电功率上限均为$5000\ \mathrm{kW}$，效率均为$0.9$。
    \item \textbf{充放电互斥}：同一时段不同时充、放电。规划层按命题1处理，执行层按命题3处理，并加入极小的充放电总量惩罚项。
    \item \textbf{信息边界}：预测、购电计划和合同调整只能使用决策时刻之前已经获得的信息，不能使用未来真实负荷、光伏或电价。
\end{enumerate}

\subsection{统一时间口径与量纲转换}
单时段为10分钟，即$\Delta t=1/6\ \mathrm{h}$。功率按$E=P\Delta t$换算为电量，负荷、光伏、合同及储能动作均以kWh计；第$t$行数据对应区间$[(t-1)\times10,t\times10)\ \mathrm{min}$。仿真和SOC更新均按这一时间口径执行。

\subsection{基础功率平衡物理模型}
记$L_t,V_t,g_t,w_t,v_t,x_t,y_t,u_t,S_t$分别为负荷、可用光伏、合同购电、闲置合同、光伏消纳、充电、放电、应急购电和时段初储电量。

时段能量平衡方程：
\[
g_t - w_t + v_t + y_t + u_t = L_t + x_t
\]

储能荷电状态递推关系：
\[
S_{t+1}=S_t+\eta_c x_t-\dfrac{y_t}{\eta_d},\quad \eta_c=\eta_d=0.9
\]

系统约束条件：
\[
\begin{cases}
0\le v_t \le V_t,\quad 0\le w_t \le g_t,\quad u_t \ge 0\\
0\le x_t,y_t \le B_{\max}\\
1200 \le S_t \le 10800
\end{cases}
\]

\subsection{模型符号说明}
{\scriptsize
\begin{longtable}{p{0.27\textwidth}p{0.13\textwidth}p{0.51\textwidth}}
\caption{模型符号说明}\label{tab:symbol}\\
\toprule
符号 & 单位 & 含义说明 \\
\midrule
\endfirsthead
\toprule
符号 & 单位 & 含义说明 \\
\midrule
\endhead
\bottomrule
\endfoot
\multicolumn{3}{l}{\textbf{决策变量}}\\
$g_t$ & kWh & 时段$t$实际生效的合同购电量。问题一中为确定性购电量。 \\
$p_t$ & kWh & 每日0:00签订的日前合同购电量。 \\
$q_t$ & kWh & 问题三中经过日内调整后最终生效的合同购电量。 \\
$\delta_t^+,\delta_t^-$ & kWh & 合同上调量和下调量，分别表示对尚未执行合同的增加量与减少量。 \\
$v_t$ & kWh & 实际用于供给负荷或储能充电的光伏电量，满足$0\le v_t\le V_t$。 \\
$x_t,y_t$ & kWh & 储能充电量和放电量。 \\
$u_t$ & kWh & 真实执行阶段为弥补供电缺口而购买的应急电量。 \\
$w_t$ & kWh & 已签合同中未被实际消纳的闲置电量，按合同价格计费。 \\
\midrule
\multicolumn{3}{l}{\textbf{参数与状态量}}\\
$L_t$ & kWh & 时段$t$的实际负荷电量。 \\
$V_t$ & kWh & 时段$t$可利用的光伏电量上限。 \\
$S_t$ & kWh & 时段$t$开始时的储能电量。 \\
$S_{\min},S_{\max}$ & kWh & 储能允许的最低、最高电量，本文分别为1200和10800。 \\
$B_{\max}$ & kWh & 单时段充、放电电量上限，$B_{\max}=P_{\max}\Delta t$。 \\
$c_t$ & 元/kWh & 时段$t$的计划购电单价，应急购电单价为$5c_t$。 \\
$\eta_c,\eta_d$ & $-$ & 储能充电效率和放电效率，本文均取0.9。 \\
$\hat L_{d,t},\hat V_{d,t}$ & kWh & 负荷和光伏的中心点预测值。 \\
$\widetilde L_{d,t},\widetilde V_{d,t}$ & kWh & 加入风险裕度后用于合同规划的保守预测值。 \\
$\hat N_{d,t},\widetilde N_{d,t}$ & kWh & 中心净负荷预测和安全净负荷预测，其中$N=L-V$。 \\
$e^L_{d,t},e^V_{d,t},e^N_{d,t}$ & kWh & 负荷、光伏和净负荷的预测残差。 \\
$Q_\alpha$ & kWh & 历史残差的$\alpha$分位数，用于构造风险裕度。 \\
$J(p)$ & 元 & 单时段合同购电与应急购电组成的期望费用函数。 \\
$A_t$ & kWh & 预测盈余，定义为合同购电加光伏预测减去负荷预测。 \\
$W$ & 日 & 计算经验分位数时使用的历史残差窗口长度。 \\
$\alpha$ & $-$ & 风险分位水平。 \\
$H$ & 时段 & 滚动LP每次向前看的时域长度。 \\
$\lambda$ & 元/kWh & 终端储能价值系数，仅用于引导优化，不计入实际费用。 \\
$\varepsilon$ & 元/kWh & 充放电总量的极小惩罚系数，用于选取数值上不同时充、放电的解。 \\
\midrule
\multicolumn{3}{l}{\textbf{集合与信息}}\\
$d,t,j,h$ & $-$ & 日、10分钟时段、历史日和预测提前期索引。 \\
$\mathcal H_t$ & $-$ & 从当前时段$t$开始的滚动优化时域。 \\
$I_{d,t}$ & $-$ & 第$d$日第$t$时段决策时已经获得的信息集。 \\
$\mathcal F_{d,h}$ & $-$ & 第$d$日提前期为$h$时可使用的历史预测信息集。 \\
\bottomrule
\end{longtable}
}

信息集定义为
\[
\begin{aligned}
I_{d,t}=\{&\text{时隙 }t\text{ 前已完成的负荷、光伏与电价观测},\ S_{d,t},\\
&\text{截至 }t\text{ 已发布的预测、合同与模型参数}\}.
\end{aligned}
\]
因此，$I_{d,t}$不包含第$d$日尚未发生时隙的真实负荷、光伏或电价。所有计划与调度变量必须是$ I_{d,t}$的函数；真实值仅在该时隙决策完成后进入执行修正与结算。
